\section{Introduction}\label{sec:intro}

The claim under examination in this paper is the strongest version of the
single-couple human origin (SHCO) hypothesis: that a founding couple of $N=2$
individuals could be the \emph{sole genetic progenitors} of a human
population --- the source of all of its heritable variation --- in the sense
that no prior population is required to explain the diversity the population
subsequently carries. The distinction between this claim and the weaker
bottleneck-survival claim is the organizing theme of the paper, and it is
worth stating at the outset in the terms in which the recent literature has
framed it. A founding couple
\emph{sampled} from a large pre-existing gene pool carries, on average, three
quarters of that pool's heterozygosity \citep{SlatkinExcoffier2012}; a
population founded by such a couple and expanded rapidly retains most of what
they carry \citep{Nei1975,BartonCharlesworth1984}. Both statements are
population-genetically routine. Neither, however, addresses the stronger
question, because both take the couple's own variation as a random draw from
an external reservoir. If the couple's genomes are instead treated as the
\emph{initial condition} of the model --- carrying whatever divergence the
biology of two primate genomes permits --- then the external reservoir
disappears from the theory, and the question splits in two. The first
question is quantitative: how much of today's diversity can be generated
below the couple, and how much must be carried into it? Section~\ref{sec:bound}
answers with an exact inequality --- mutation since $T_0$ can supply at most
$2\mu T_0 \approx 53\%$ of the observed $\pi$, so an admissible couple must
carry $h \ge \tfrac{4}{3}(\pi_{\mathrm{obs}} - 2\mu T_0) \approx 0.63\times10^{-3}$,
some $1.9$ million SNVs per founder --- and the canonical endowment clears
this floor comfortably. The second question is provenance, and it must be
addressed before the model is stated, because the objection it invites is
raised at the first equation. The canonical endowment,
$h = 1.576\times10^{-3}$, is the coalescent-equivalent of $2\Ne$ at
$\Ne \approx 3.3\times10^{4}$: numerically the $\Nepre$ of the
sampled-couple class. Has the pre-bottleneck population therefore merely
been renamed? It has not, and the reason is worth stating at the outset,
since everything downstream depends on it. What the two classes share is a
single scalar --- a depth reservoir that any origin must supply somewhere,
because mutation alone accounts for at most $53\%$ of the observed
diversity. What the sampled class adds to that scalar, and the endowed-pair
class does not, is an \emph{entity}: a population, with drift, structure,
and a genealogy reaching above the couple, whose presence is not one more
parameter but a source of signatures --- pre-split allelic lineages at
balanced loci, deep structured coalescences, deme-private variants --- that
diversity measurements are in principle able to detect. The endowed-pair
class contains no such entity: above the couple the model's genealogy is
deliberately undefined, no pre-split lineage exists anywhere in it, and the
entire content of the assumption is one bounded scalar, its admissible
values fixed by the heterozygosities of observed primate genomes and its
provenance deliberately left open. Three readings of that scalar are live
--- sampled from a pre-population, constituted de novo, assembled from
diverged ancestral lineages --- and all three produce identical descendant
genetics (Section~\ref{sec:nature}), which is precisely why no diversity
measurement can adjudicate among them and why a theory that claims to is
importing an assumption, not deriving a result. What the model establishes
is underdetermination under exact bounds: the historical diversity
objection, which purported to read the couple's impossibility off the data,
follows from no diversity measurement.

The distinction is not academic. It maps onto the modern debate over whether
genomic data constrain the existence of a single founding couple --- a debate
carried across population genetics, philosophy of science, and the
theological literature \citep{Gauger2012,Swamidass2019,HossjerGauger2019}.
The standard objection --- that humans
are far too diverse to descend from a single pair --- has two halves: a
\emph{level} claim (a bottleneck to $N=2$ would leave the species nearly
homozygous) and a \emph{structure} claim (the allele-frequency spectrum, the
coalescence-time distribution, and the allelic richness of balanced loci
would betray such an origin). The level claim fails on the classical results:
a one-generation bottleneck at $N=2$ costs exactly one quarter of the
carried heterozygosity, and rapid expansion preserves the remainder
\citep{Nei1975,SlatkinExcoffier2012}. To this the paper adds the decisive
combinatorial point: because the overwhelming majority of human SNPs are
biallelic, four gene copies can carry \emph{every} allelic state at every
biallelic locus simultaneously, so the carrying capacity of two genomes is
not the binding constraint. The structure claims are the
substantive ones, and they have never been jointly tested against a
sole-progenitor model: the sampled-couple class tested against them
turns out to be a carrier model by construction --- it assumes the
pre-bottleneck population whose existence the sole-progenitor claim denies.
The present paper closes that gap.

The strategy is to work entirely within secular
population genetics, state every constraint as an exact identity or a
published empirical estimate, and map the parameter region in which all
constraints hold simultaneously. Nothing in the argument requires
nonstandard genetics: every rate is pedigree-measured, every correction is
a published reconciliation, every bound is an observed genome, and the
model's one free parameter is fixed within narrow margins by the
heterozygosities of genomes that exist. The
\emph{endowed-pair (EP) model} introduced here has no pre-bottleneck
population at all. Its founder genomes enter with an endowment $h$ --- the
per-site divergence among their four gene copies --- bounded above not by a
population history but by the heterozygosity that natural primate genomes
are observed to carry: $3.1$--$4.1$ million SNVs in humans
\citep{PradoMartinez2013,Auton2015}, $5.7$ million in the chimpanzee
\citep{PradoMartinez2013}. The endowment is worth $h/(2\mu)$ generations of
coalescent-equivalent depth --- $1.64$ Myr at the canonical value, $1.98$ Myr
for chimp-scale founders --- which is precisely the depth that the
sampled-couple model imported from its $\Nepre$. The master result (Section~\ref{sec:endo})
is that the pairwise-diversity identity $\pi = 2\mu E[T_2]$ is satisfied at
the observed value $\pi \approx 10^{-3}$ by a region of parameter space
whose bottleneck time coincides with the window that the independent
autosomal arguments of Swamidass already singled out: $T_0 \approx 400$
to $600$ kya, with the compression time of four lineages --- the quantity
Swamidass named TMR4A and estimated at $495 \pm 100$ kya
\citep{Swamidass2019,Swamidass2022} --- equal to $T_0$ identically.

Three further results complete the feasibility map. First, the
site-frequency spectrum, Tajima's $D$, and the linkage-disequilibrium decay
of the EP model are computed by forward simulation and coalescent
simulation; the model reproduces the strongly negative $D$ of expanding
human populations and the observed short-range $r^2$ profile
(Section~\ref{sec:sfsld}). Second, the extreme bottleneck is invisible to
PSMC- and MSMC-type reconstruction at its temporal resolution
(Section~\ref{sec:ne}), a conclusion that follows from the resolution limits
of the reconstruction methods \citep{LiDurbin2011} and is confirmed by the
present simulations. Third, the
polymorphism objection at balanced loci is answered with a retention and
regeneration calculus, supplemented by the
convergent-evolution and bottleneck-duration corrections that the MHC
literature itself supplies against the classical HLA argument of
\citet{Ayala1995}
(Section~\ref{sec:poly}). Fourth, the introgressed archaic fraction of
modern genomes --- some $2$--$6\%$ of non-African DNA --- is handled by the
descendant-clade extension (Section~\ref{sec:scope}): the Neanderthal and
Denisovan lineages are modeled as descendant branches of the founding
couple's descendant population, so that archaic admixture is gene flow
\emph{within} the couple's pedigree and the introgressed tracts are
couple-descended DNA returning to the modern branch.

The paper is equally explicit about what the model does \emph{not} settle.
Under sole progenitorship the mitochondrial-DNA and Y-chromosomal most
recent common ancestors are the founding couple themselves, so their true
dates equal $T_0$ identically; their published date estimates, however, are
anchored on the human--chimpanzee split, a clock that \citet{Scally2012}
showed requires an ancestral mutation rate roughly twice the directly
measured trio rate. Section~\ref{sec:uni} quantifies the correction and
identifies it as the model's principal conditional layer. And the
introgressed archaic fraction of modern genomes is subject to the
descendant-clade extension and its scoping corollary
(Section~\ref{sec:scope}), whose clock-free tests delimit exactly which
versions of the sole-progenitor claim the genetic data leave open. The
claim defended here is a possibility theorem: the diversity and structure of
the human genome do not falsify a sole-progenitor couple at
$T_0 \in [400, 600]$ kya --- and the theorem is the more falsifiable of
the two positions on the question, because Section~\ref{sec:theorem}
states seven measurable conditions under which it collapses, while the
counter-position, that the couple must have been sampled from a
pre-existing population, rests on an assumption that the
identical-observables result places beyond the reach of any diversity
measurement (Section~\ref{sec:nature}). Whether the claim's metaphysical
provenance is creation, sampling, or assembly from diverged lineages is
not a question the mathematics can decide, and the paper does not pretend
otherwise; what the mathematics does decide is that no measurement of
diversity adjudicates it --- which is why the historical objection, posed
as a claim about the data, never followed from the data.

The plan of the paper is as follows. Section~\ref{sec:framework} fixes
notation and the coalescent background. Section~\ref{sec:classes} defines
the two model classes and the endowment parameter with its biological
bounds. Section~\ref{sec:endo} derives the master equation of the EP model,
prices the generation bound the endowment must meet, and maps its
feasibility region. Section~\ref{sec:uni} treats the
uniparental clocks; Section~\ref{sec:sfsld} the frequency spectrum,
Tajima's $D$, linkage disequilibrium, and coalescence-depth distribution;
Section~\ref{sec:ne} the detectability of the bottleneck;
Section~\ref{sec:poly} the balanced-locus objection;
Section~\ref{sec:scope} the descendant-clade extension of the archaic
fraction and its scoping corollary; Section~\ref{sec:emp} the
theoretical and empirical existence proofs. Section~\ref{sec:theorem}
assembles the feasibility theorem and its falsification conditions,
Section~\ref{sec:objections} answers the objection catalogue in turn, and
Sections~\ref{sec:nature}--\ref{sec:conclusion} discuss the proof's nature,
scope, and conclusions. All quantitative statements are emitted by a single
fixed-seed script (Section~\ref{sec:sims} and the reproduction appendix).

\section{Mathematical framework}\label{sec:framework}

\subsection{The Wright--Fisher process and the coalescent}
The backward genealogy of the Wright--Fisher process \citep{Wright1931}
of a sample from a Wright--Fisher population of changing diploid size
$N(t)$ is governed by the pairwise coalescence hazard
\begin{equation}\label{eq:hazard}
\Lambda(t) \;=\; \int_0^{t} \frac{\mathrm{d}s}{2\,N(s)} ,
\end{equation}
the probability that a pair of gene copies survives un-coalesced from the
present to backward time $t$ being $e^{-\Lambda(t)}$
\citep{Hudson1990,CrowKimura1970}. The identity that governs the diversity
analysis is the infinite-sites relation
\begin{equation}\label{eq:piconstraint}
\pi \;=\; 2\mu\,E[T_2],
\end{equation}
exact for neutral mutations on any genealogy and under any demography
\citep{Hudson1990}. With the observed genome-wide nucleotide diversity
$\pi_{\mathrm{obs}} \approx 10^{-3}$ per site \citep{Auton2015} and the
directly measured per-generation mutation rate
$\mu = 1.2\times10^{-8}$ \citep{Kong2012,Scally2012,Jonsson2017}, the
requirement is
\begin{equation}\label{eq:T2req}
E[T_2] \;=\; \frac{\pi_{\mathrm{obs}}}{2\mu} \;\approx\; 4.17\times10^{4}\ \text{generations}
\;\approx\; 1.04\ \text{Myr}
\end{equation}
at 25 years per generation \citep{Fenner2005,Scally2012}. Equation~\eqref{eq:T2req} is the master
constraint of the paper; what
changes between model classes is the mechanism by which an $N=2$ origin meets it.

Two clock classes must be distinguished from the outset, because the
feasibility argument of Section~\ref{sec:uni} depends on the distinction.
\emph{Direct} clocks are calibrated on pedigrees: the autosomal trio rate
$\mu = 1.0$--$1.3\times10^{-8}$ \citep{Scally2012,Jonsson2017}, the
recombination map, and every autosomal coalescence quantity derived from
them. \emph{Anchored} clocks are calibrated on the human--chimpanzee
divergence: the mtDNA and Y-chromosomal rates and their coalescence-date
estimates. The two classes disagree systematically in the literature
precisely because the anchored clocks embed an ancestral rate roughly twice
the direct one \citep{Scally2012}; Section~\ref{sec:uni} makes this
quantitative. Autosomal quantities ($\pi$, $E[T_2]$, TMR4A, LD) are read on
the direct clock throughout.

\subsection{Summary statistics}
Heterozygosity per site $H$, nucleotide diversity $\pi$, Watterson's
$\theta_W = S_n/a_n$ with $a_n = \sum_{i=1}^{n-1} 1/i$
\citep{Watterson1975}, Tajima's $D$ \citep{Tajima1989}, and the linkage
disequilibrium coefficient $r^2$ \citep{HillWeir1988} are used with their
standard definitions and simulation-based expectations. Tajima's $D$ for the
human series is strongly negative
($D \approx -2$ in large-sample series \citep{Tennessen2012}), the
signature of recent expansion; the EP model is required to reproduce its
sign and order of magnitude (Section~\ref{sec:sfsld}).

\subsection{Notation}
Table~\ref{tab:notation} summarizes notation; one generation is 25 years
\citep{Fenner2005,Scally2012}. All times are generations unless marked ``kya''.

\begin{table}[t]
\centering\small
\caption{Notation used throughout.}
\label{tab:notation}
\begin{tabular}{ll}
\toprule
Symbol & Meaning \\
\midrule
$N(t)$, $\Ne^{\max}$ & diploid effective size at backward time $t$; post-expansion plateau \\
$T_0$ & time of the founding couple (generations before present) \\
$\gamma$ & per-generation growth factor of the early expansion \\
$h$ & endowment: per-site divergence among the couple's four gene copies \\
$\Lambda_c$ & hazard accumulated from present to the couple generation, Eq.~\eqref{eq:hazard} \\
$D_{<c}$ & expected coalescence depth contributed by pairs coalescing before the couple \\
$\mu$, $c$ & per-site per-generation mutation ($1.2\times10^{-8}$) and recombination ($10^{-8}$) rates \\
$H$, $\pi$ & heterozygosity per site; nucleotide diversity \\
$T_2$, $\mathrm{TMR4A}$ & pairwise coalescence time; four-lineage genome compression time \\
$F$ & inbreeding coefficient \\
\bottomrule
\end{tabular}
\end{table}

\section{The two model classes}\label{sec:classes}

\subsection{Sampled-couple trajectories}
The \emph{sampled-couple trajectory} $B(T_0, \Nepre)$ is the baseline model
class against which the sole-progenitor question is usually posed: a
demographic function with $N(T_0) = 2$, the
founding pair being a random sample of four gene copies from a pre-existing
population of effective size $\Nepre$. The pass-through structure of that
model imports the pre-population's depth: pairs that survive the expansion
coalesce, on average, in a gene pool of size $\Nepre$, giving
\begin{equation}\label{eq:sampled}
E[T_2] \;=\; D_{<c} \;+\; e^{-\Lambda_c}\Big[T_0 + \tfrac{3}{4}\,(2\Nepre)\Big],
\end{equation}
so that observed diversity forces $\Nepre \gtrsim 2$--$3\times10^4$. This
is the quantitatively correct statement of the ``carrier'' result: in the
sampled-couple class, the couple is a bottleneck transmitter, and the
pre-bottleneck population is a mandatory input. A feasibility region for
that class follows from the same hazard machinery developed in
Section~\ref{sec:endo}; the present paper does not re-derive it, because the
class is by construction a carrier model --- and it is precisely the
carrier assumption that sole progenitorship denies.

\subsection{Endowed-pair trajectories}
The \emph{endowed-pair (EP) trajectory} $E(T_0, \gamma, \Ne^{\max}, h)$
removes the pre-bottleneck population. Its demographic function has
$N(T_0) = 2$ --- the founding couple carrying four gene copies per locus ---
and no population before the couple: the couple is the model's first
generation, expanding forward at factor $\gamma$ per generation to a
plateau $\Ne^{\max}$, with the recent-growth epoch of the last two
millennia appended. What replaces $\Nepre$ is an
initial condition on the couple's genomes:

\begin{definition}[Endowment]\label{def:endowment}
The \emph{endowment} of a founding couple is the per-site divergence $h$
among its four gene copies at time $T_0$: the fraction of sites at which
the couple's homologous copies differ. The corresponding pool heterozygosity
--- the expected heterozygosity of two gene copies drawn at random from the
couple's four --- is $H_{\mathrm{pool}} = \tfrac{3}{4}h$ for founders with
disjoint SNV sets and comparable heterozygosity (Section~\ref{sec:endocap}).
\end{definition}

Three remarks fix the interpretation. First, the endowment is a free
parameter of the model class, not a derived quantity: under the
sole-progenitor reading there is no ancestral pool from which to sample it.
Second, the parameter is bounded \emph{biologically}, not historically: two
primate genomes cannot carry more divergence than natural primate genomes
in fact carry (Section~\ref{sec:endocap}). Third, the model's genealogy
above the couple is deliberately undefined --- cross-founder lineages have
no common ancestor within the model, and their divergence is the endowment
itself. This is the formal expression of sole progenitorship: nothing
before the couple enters the calculation.

The EP model is the mathematical skeleton of the scenario in which a
bottleneck of two lies beyond half a million years ago, and
it is the model class in which the couple is a \emph{source} rather than a
carrier. The two classes are complementary, not competing: both are
consistent with the same data (Sections~\ref{sec:endo}--\ref{sec:sfsld}),
and they differ in where the diversity enters --- externally through
$\Nepre$, or internally through $h$. Which class describes actual history
is not decidable by the diversity data alone; the difference is one
hierarchical level of ancestry (Section~\ref{sec:nature}).

\subsection{Biological bounds on the endowment}\label{sec:endocap}
The endowment's admissible range is fixed by the heterozygosity that
natural primate genomes carry. The average modern human genome carries
about $3.1$--$4.1$ million single-nucleotide variants
\citep{PradoMartinez2013,Auton2015}; the chimpanzee, at comparable genome
size, carries about $5.7$ million \citep{PradoMartinez2013}. Converting to
per-site divergences at $L = 3\times10^{9}$ sites:
\begin{equation}\label{eq:hbounds}
h_{\mathrm{human}} = 1.36\times10^{-3}, \qquad
h_{\mathrm{chimp}} = 1.90\times10^{-3},
\end{equation}
with the corresponding pool heterozygosities $0.75h$ of
$1.02\times10^{-3}$ and $1.43\times10^{-3}$. Two founders whose variant
sets are largely disjoint and whose heterozygosities lie anywhere in the
observed primate range therefore endow the model with a pool heterozygosity
of order $10^{-3}$ --- the entire diversity budget that
Eq.~\eqref{eq:T2req} demands. The carrying-capacity objection dissolves at
this step: the couple need not carry ``all'' of today's $10^8$-scale SNP
catalog, because the rare-variant tail of that catalog is generated by
post-bottleneck mutation (Section~\ref{sec:sfsld}); what the couple must
carry is the \emph{divergence budget} of the common-variant backbone, and
that budget is of the same order as two ordinary primate genomes.

The combinatorial version of the same point is elementary:
at a biallelic locus the couple's four gene copies can carry every allelic
state that exists, and sexual reproduction reassorts them independently;
the four-allele endowment is therefore representable in full. The
transmission mechanics are quantified in Section~\ref{sec:trans}.
